\section{Independent Mass Prompt Confirmation}
\label{app:mass-confirmation}

\subsection{Cohorts, prompt factors, and calibration}

The prospective confirmation fixes Qwen2.5-VL-7B, the consumed
\texttt{model.visual.blocks.31}, $336\times336$ images, the six train-fitted
clinical normals, and relative-token dose $\alpha=+0.25$.
Its 200 patients are absent from the initial intervention, Effusion
specificity, causal-ownership, and Consolidation closure cohorts.
Selection uses only patient, row, and split metadata: sorted eligible
patient IDs are permuted with seed 20260909, the first 200 are retained,
and the same generator selects one sorted row per patient.
The crossover's 200 calibration patients are reused in their stored order
and remain disjoint from confirmation. Calibration contains 15 positive
and 185 negative patients; confirmation contains 12 positive and 188 negative
patients.

Figure~\ref{fig:mass-prompts} shows the prompt factors. Each of the two
question wordings is concatenated with each of the three explicit
instructions, yielding six cells. A seventh anchor reproduces the original
yes/no prompt. Thus the two \textit{show} coded cells reproduce the
discovery prompts exactly, while the factorial comparisons vary one factor
at a time.

\begin{figure}[H]
\centering
\begin{minipage}{0.95\linewidth}
\small
\fbox{\parbox{0.95\linewidth}{Does this chest radiograph show a lung mass?\par
Is there a lung mass in this chest radiograph?}}
\par\smallskip\centerline{(a) Question wordings}\medskip
\fbox{\parbox{0.95\linewidth}{Answer yes if the finding is present and no if it is absent. Reply with yes or no only.\par\smallskip
Answer A if the finding is present and B if it is absent. Reply with A or B only.\par\smallskip
Answer B if the finding is present and A if it is absent. Reply with A or B only.}}
\par\smallskip\centerline{(b) Explicit response instructions}\medskip
\fbox{\parbox{0.95\linewidth}{Is there a lung mass in this chest radiograph? Answer yes or no.}}
\par\smallskip\centerline{(c) Original-prompt anchor}
\end{minipage}
\caption{\textbf{The Mass confirmation crosses question wording with response instructions.} Each question in (a) combines with each instruction in (b); (c) supplies the seventh condition.}
\label{fig:mass-prompts}
\end{figure}

For each patient and cell, the finding-present logit margin is yes-minus-no,
A-minus-B, or B-minus-A according to the response instruction.
Scores use the maximum logit over accepted single-token spellings, including
leading-space variants and the accepted yes/no case variants.
Each intervention effect subtracts that cell's own unsteered margin.
The seven cells each receive a baseline, six clinical directions,
119 shared seed-0 isotropic random directions, and the Mass
coordinate-permutation sham drawn after the random family. This yields
177,800 confirmation and 1,400 calibration outcomes. The complete run
took 8,783 seconds on one A100.

A cell is eligible when calibration contains at least ten patients per
class, at least 1,900 of 2,000 bootstrap draws have both classes, and the
one-sided 95\% AUROC lower bound exceeds 0.5. Calibration resampling uses
seed 20260908. Table~\ref{tab:mass-calibration} reports the clean answer
discrimination and Brier error for all seven cells.

\begin{table}[H]
\centering
\caption{Mass capability calibration across prompt conditions.}
\label{tab:mass-calibration}
\resizebox{\ifdim\width>\linewidth\linewidth\else\width\fi}{!}{%
\begin{tabular}{lrrr}
\toprule
Condition & AUROC & One-sided 95\% lower bound & Brier score \\
\midrule
\textit{show}/yes/no & 0.6486 & 0.5105 & 0.0703 \\
\textit{show}/A-present & 0.6533 & 0.5198 & 0.0894 \\
\textit{show}/B-present & 0.6542 & 0.5231 & 0.0939 \\
\textit{is}/yes/no & 0.6422 & 0.5043 & 0.0720 \\
\textit{is}/A-present & 0.6420 & 0.5047 & 0.1006 \\
\textit{is}/B-present & 0.6436 & 0.5051 & 0.1041 \\
anchor & 0.6443 & 0.5088 & 0.0716 \\
\bottomrule
\end{tabular}%
}
\end{table}

All cells are eligible, with 2,000 valid draws each. Their AUROCs range
from 0.6420 to 0.6542 and their lower bounds from 0.5043 to 0.5231.
Eligibility is determined on calibration before inspecting intervention
outcomes; the confirmation cohort supplies a separate assessment of
predictive behavior below.

\subsection{Clinical competition and the registered control family}

The primary family comprises the four coded cells' five Mass-minus-other
clinical contrasts, for 20 comparisons in total. Joint resampling of the
200 confirmation patients preserves pairing across cells and directions.
A centered studentized max-$T$ bootstrap with 5,000 draws and seed 20260910
gives simultaneous one-sided 95\% lower bounds using
Equation~\ref{eq:max-t}: each contrast's bootstrap standard deviation
is computed once across the draws, with denominator 4,999, and held
fixed during standardization. The degenerate-component rule is given in
Appendix~\ref{app:simultaneous-bounds}. A coded cell qualifies when
it is eligible, all five bounds are positive, and its positive Mass effect
exceeds every same-cell random effect and the absolute sham effect.
Cross-wording specificity requires all four coded cells to qualify;
original-wording replication requires both \textit{show} coded cells under
the same joint family. The explicit yes/no cells and anchor are descriptive.

Table~\ref{tab:mass-controls} reports the seven mean effects and clinical
margins. A clinical margin subtracts the largest of the five competing
clinical means within the cell. Its descriptive 95\% percentile interval
recomputes that maximum in each paired-patient bootstrap draw; it is
distinct from the simultaneous pairwise bounds in
Table~\ref{tab:mass-confirmation}.

\begin{table}[H]
\centering
\caption{Mass clinical and generic-control competition.}
\label{tab:mass-controls}
\resizebox{\ifdim\width>\linewidth\linewidth\else\width\fi}{!}{%
\begin{tabular}{lrrrrr}
\toprule
Condition & Mass effect & Clinical margin & 95\% interval & Random max & $|$Sham$|$ \\
\midrule
\textit{show}/yes/no & 2.3800 & $0.2369$ & $[0.1362,0.3394]$ & 3.0531 & 0.1225 \\
\textit{show}/A-present & 1.7019 & $0.2525$ & $[0.1869,0.3175]$ & 2.2100 & 0.2931 \\
\textit{show}/B-present & 1.8188 & $0.2338$ & $[0.1675,0.3006]$ & 2.2169 & 0.2850 \\
\textit{is}/yes/no & 2.6600 & $-0.4556$ & $[-0.5350,-0.3794]$ & 3.6031 & 0.6556 \\
\textit{is}/A-present & 1.9225 & $-0.1087$ & $[-0.1606,-0.0575]$ & 2.4750 & 0.6075 \\
\textit{is}/B-present & 2.0019 & $-0.2006$ & $[-0.2531,-0.1487]$ & 2.4944 & 0.5944 \\
anchor & 2.9987 & $-0.7438$ & $[-0.8338,-0.6575]$ & 4.1312 & 0.9856 \\
\bottomrule
\end{tabular}%
}
\end{table}

Effusion is the strongest clinical competitor in every cell.
The \textit{show} A-present and B-present clinical margins are positive,
with minimum simultaneous lower bounds 0.1675 and 0.1462, respectively.
Their Mass effects of 1.7019 and 1.8188 remain below random maxima 2.2100
and 2.2169. Two of the 119 random effects exceed Mass in each cell,
giving reference ranks of 3/120; the criterion requires exceeding all 119.
Under \textit{is}, the two coded clinical margins are negative.
Consequently, both registered specificity conjunctions are unmet.
Clinical competition reproduces under the discovery wording, while the
larger random family limits the supported specificity claim.

\subsection{Paired wording and instruction contrasts}

Table~\ref{tab:mass-contrasts} isolates the observed changes in clinical
competition. Each row subtracts two margins on the same patients, with
the clinical maximum recomputed within each cell and resample.
These eight paired percentile intervals are descriptive and do not change
the two registered conjunctions.

\begin{table}[H]
\centering
\caption{Paired changes in Mass clinical competition.}
\label{tab:mass-contrasts}
\resizebox{\ifdim\width>\linewidth\linewidth\else\width\fi}{!}{%
\begin{tabular}{lrr}
\toprule
Clinical-margin contrast & Difference & 95\% interval \\
\midrule
\textit{show}/A-present $-$ \textit{show}/yes/no & $0.0156$ & $[-0.0250,0.0544]$ \\
\textit{show}/B-present $-$ \textit{show}/yes/no & $-0.0031$ & $[-0.0425,0.0350]$ \\
\textit{is}/A-present $-$ \textit{is}/yes/no & $0.3469$ & $[0.3150,0.3788]$ \\
\textit{is}/B-present $-$ \textit{is}/yes/no & $0.2550$ & $[0.2219,0.2881]$ \\
\textit{is}/yes/no $-$ \textit{show}/yes/no & $-0.6925$ & $[-0.7369,-0.6494]$ \\
\textit{is}/A-present $-$ \textit{show}/A-present & $-0.3612$ & $[-0.3906,-0.3319]$ \\
\textit{is}/B-present $-$ \textit{show}/B-present & $-0.4344$ & $[-0.4644,-0.4056]$ \\
\textit{is}/yes/no $-$ anchor & $0.2881$ & $[0.2631,0.3137]$ \\
\bottomrule
\end{tabular}%
}
\end{table}

Under \textit{show}, the coded-minus-explicit-yes/no intervals span zero;
they leave the instruction contrast unresolved rather than establishing
equivalence. Under \textit{is}, both coded instructions improve the margin
relative to explicit yes/no, although their margins remain negative.
Within every response instruction, changing \textit{show} to \textit{is} lowers the
clinical margin. Explicit yes/no instructions also improve the
\textit{is} margin relative to the original anchor.
The joint pattern distinguishes wording-conditioned clinical competition
from a claim that A/B symbols alone create specificity.

\subsection{Answer discrimination and probability error}

Table~\ref{tab:mass-prediction} compares baseline and Mass-steered answer
AUROC on the 200 confirmation patients. Scores are clinically oriented
answer margins, as in capability calibration; they are not fitted probe
scores. The paired change intervals use the same 5,000 patient resamples.

\begin{table}[H]
\centering
\caption{Mass steering and answer discrimination.}
\label{tab:mass-prediction}
\resizebox{\ifdim\width>\linewidth\linewidth\else\width\fi}{!}{%
\begin{tabular}{lrrrr}
\toprule
Condition & Baseline AUROC & Steered AUROC & Change & 95\% interval \\
\midrule
\textit{show}/yes/no & 0.4938 & 0.5288 & $0.0350$ & $[-0.1505,0.2129]$ \\
\textit{show}/A-present & 0.5055 & 0.5461 & $0.0406$ & $[-0.1347,0.2136]$ \\
\textit{show}/B-present & 0.4962 & 0.5468 & $0.0505$ & $[-0.1307,0.2223]$ \\
\textit{is}/yes/no & 0.4942 & 0.4989 & $0.0047$ & $[-0.1728,0.1735]$ \\
\textit{is}/A-present & 0.4951 & 0.5202 & $0.0250$ & $[-0.1467,0.1911]$ \\
\textit{is}/B-present & 0.4940 & 0.5441 & $0.0501$ & $[-0.1451,0.2263]$ \\
anchor & 0.4956 & 0.4807 & $-0.0148$ & $[-0.1802,0.1435]$ \\
\bottomrule
\end{tabular}%
}
\end{table}

Baseline AUROCs lie near 0.5 in this cohort, and all seven paired
AUROC-change intervals span zero. With only 12 positive confirmation
patients, the intervals are broad. The calibration-based eligibility
decision and this imprecise confirmation-cohort discrimination estimate
answer different questions.

Table~\ref{tab:mass-brier} evaluates probability error, using the sigmoid
of each finding-present logit margin. The Brier score is the mean squared
difference between that probability and the binary disease label;
positive changes indicate larger error.

\begin{table}[H]
\centering
\caption{Mass steering and probability error.}
\label{tab:mass-brier}
\resizebox{\ifdim\width>\linewidth\linewidth\else\width\fi}{!}{%
\begin{tabular}{lrrrr}
\toprule
Condition & Baseline Brier & Steered Brier & Change & 95\% interval \\
\midrule
\textit{show}/yes/no & 0.0601 & 0.1530 & 0.0929 & $[0.0691,0.1163]$ \\
\textit{show}/A-present & 0.0858 & 0.2949 & 0.2091 & $[0.1829,0.2338]$ \\
\textit{show}/B-present & 0.0943 & 0.3131 & 0.2188 & $[0.1911,0.2445]$ \\
\textit{is}/yes/no & 0.0644 & 0.2238 & 0.1594 & $[0.1305,0.1871]$ \\
\textit{is}/A-present & 0.0985 & 0.3866 & 0.2881 & $[0.2589,0.3153]$ \\
\textit{is}/B-present & 0.1067 & 0.4015 & 0.2948 & $[0.2644,0.3227]$ \\
anchor & 0.0630 & 0.2189 & 0.1559 & $[0.1269,0.1836]$ \\
\bottomrule
\end{tabular}%
}
\end{table}

Mass steering raises finding-present probabilities while increasing Brier
error in every cell, with all seven change intervals above zero.
Thus a positive clinical competition margin under \textit{show} coexists
with worse probability predictions and unresolved discrimination change.
These measurements concern one finding, model, visual locus, and dose;
the wording contrasts characterize that intervention interface.

\FloatBarrier
